% chapters/ch8_exp_ablation_shuffle.tex --- 4.5 アブレーション検証実験：時系列構造の寄与
\section{アブレーション検証実験：時系列構造の寄与}
\label{sec:exp_ablation_shuffle}

\subsection{実験の目的}

不確実度$U$の時系列構造が有効に機能しているかを，シャッフル実験で検証する．

\subsection{実験結果}

$U$をシャッフルすると，share100が0.593から0.943へ増加し，固定\SI{100}{\milli\second}に近づいた．省電力効果が失われることから，時系列構造が適応制御に重要であることを示した．

\clearpage
